% =====================================================================
%  10  分離戦略 — 左=分離の考え方（タイトル付き説明＋沸点マップ）｜右=工程フロー＋膜有無比較
%  source_memo: trial #201・07章 Table7.1/Fig7.1。分離目標(要求純度)を主役に:
%   C3H6 99.5wt%(α≈1.07)・H2 99.9mol%。脱ブタン=C4除去、脱エタン=軽質分離。
%   膜有無比較(共通フィード C3H6 60.7mol%): 膜+塔 TAC135/117段/94m vs 塔単独 TAC1625/200段/156m。
%  方針：まず左で「沸点→分離操作」の考え方を示し、右で工程フローと膜の効果へ。
%        左右は全高の仕切りで分ける。重なり・はみ出し・無駄余白なし。丸括弧禁止。
% =====================================================================
\definecolor{SepGreen}{HTML}{2E9E5B}
\begin{frame}[t]

  \vspace{0.05em}

  \begin{columns}[T,onlytextwidth]
    % ===== 左：分離の考え方（タイトル付き説明＋沸点マップ）｜右端に全高の仕切り =====
    \begin{column}{0.34\textwidth}
      \resizebox{!}{\dimexpr\textheight-6mm\relax}{%
      \begin{tikzpicture}[
        x=1cm, y=1cm, >=Stealth, font=\small,
        cut/.style={dashed, HeaderBlue, line width=0.8pt, -{Stealth[length=2mm]}},
        cuthard/.style={dashed, SpRed, line width=1.1pt, -{Stealth[length=2.2mm]}},
        tick/.style={HeaderBlue, line width=0.8pt},
        lab/.style={anchor=west, font=\small\bfseries},
        cmp/.style={anchor=west, font=\small},
        val/.style={anchor=east, font=\footnotesize, text=black!55},
      ]
        \useasboundingbox (-0.7,-0.9) rectangle (4.55,10.7);
        % --- 全高の仕切り縦線（右端：右の工程と分ける）---
        \draw[black!35, line width=0.9pt] (4.35,-0.8) -- (4.35,10.6);
        % --- タイトル付き説明 ---
        \node[anchor=north west, font=\large\bfseries, text=HeaderBlue] at (-0.6,10.5) {分離の考え方};
        \node[anchor=north west, font=\normalsize, text width=4.6cm, align=left] at (-0.6,9.7)
          {成分を沸点順に並べ、各境界に分離操作を割り当てた。沸点差の大きい境界は蒸留、沸点差 $6$ K と小さい C3 対のみ膜＋蒸留で分離。};
        % --- 沸点マップ（軸＝上が高沸点）---
        \draw[HeaderBlue, line width=1.3pt, ->] (0.3,0.5) -- (0.3,7.15);
        \node[anchor=west, font=\footnotesize, text=black!60] at (0.5,6.98) {[$^\circ$C]};
        \foreach \y/\name/\bp in {6.7/{$\mathrm{C_4H_{10}}$}/{$-1$}, 3.9/{$\mathrm{C_2H_6}$}/{$-89$}, 2.6/{$\mathrm{CH_4}$}/{$-161$}, 1.0/{$\mathrm{H_2}$}/{$-253$}}{
          \draw[tick] (0.15,\y) -- (0.45,\y);
          \node[cmp] at (0.5,\y) {\name};
          \node[val] at (0.0,\y) {\bp};
        }
        % C3 ペア（赤・近接）
        \draw[tick, SpRed] (0.15,5.5) -- (0.45,5.5);
        \draw[tick, SpRed] (0.15,5.1) -- (0.45,5.1);
        \node[cmp, text=SpRed] at (0.5,5.5) {$\mathrm{C_3H_8}$};
        \node[cmp, text=SpRed] at (0.5,5.1) {$\mathrm{C_3H_6}$};
        \node[val, text=SpRed] at (0.0,5.5) {$-42$};
        \node[val, text=SpRed] at (0.0,5.1) {$-48$};
        % cuts（境界ごとの分離操作）
        \draw[cut] (0.3,6.1) -- (1.9,6.1);
        \node[lab] at (2.0,6.1) {脱ブタン};
        \draw[cuthard] (0.3,5.3) -- (1.9,5.3);
        \node[lab, text=SpRed] at (2.0,5.3) {膜＋蒸留};
        \node[anchor=west, font=\footnotesize, text=SpRed] at (2.0,4.9) {沸点差 6 K};
        \draw[cut] (0.3,4.5) -- (1.9,4.5);
        \node[lab] at (2.0,4.5) {脱エタン};
        \draw[cut] (0.3,1.8) -- (1.9,1.8);
        \node[lab] at (2.0,1.8) {PSA};
        % --- マップのタイトル（最下部）---
        \node[anchor=north, font=\large\bfseries] at (1.8,-0.15) {沸点と分離操作};
      \end{tikzpicture}}
    \end{column}

    % ===== 右：工程フロー（上）＋膜有無の比較（下）=====
    \begin{column}{0.63\textwidth}
      % --- 上：横向き工程フロー（分離ブロックを緑・大きく強調）---
      \resizebox{\linewidth}{!}{%
      \begin{tikzpicture}[
        x=1cm, y=1cm, >=Stealth, font=\small,
        sep/.style={draw=SepGreen, line width=1.4pt, fill=SepGreen!22, rounded corners=3pt, align=center,
                    inner sep=2pt, minimum width=1.75cm, minimum height=1.05cm, font=\normalsize\bfseries},
        rxn/.style={draw=black, fill=white, rounded corners=2pt, align=center,
                    inner sep=2pt, minimum width=1.3cm, minimum height=0.72cm, font=\small\bfseries},
        raw/.style={align=center, font=\small\bfseries},
        goalC/.style={draw=SpRed, line width=1.2pt, fill=SpRed!12, rounded corners=3pt,
                      align=center, inner sep=4pt, text=SpRed},
        goalH/.style={draw=HeaderBlue, line width=1.2pt, fill=HiBlue!55, rounded corners=3pt,
                      align=center, inner sep=4pt},
        sub/.style={align=center, font=\scriptsize, text=black!60},
        main/.style={-{Stealth[length=2.8mm]}, HeaderBlue, line width=1.6pt},
        br/.style={-{Stealth[length=2.2mm]}, black!70, line width=1.1pt},
      ]
        % 主フロー（分離=緑で大きく、反応器=小さめ青）
        \node[raw] (raw) at (0.0,3.0) {原料\\LPG};
        \node[sep] (deb) at (2.0,3.0) {脱ブタン};
        \node[rxn] (rx)  at (4.0,3.0) {反応器};
        \node[sep] (de)  at (6.0,3.0) {脱エタン};
        \node[sep] (mc)  at (8.1,3.0) {膜＋蒸留};
        \node[goalC] (prod) at (10.8,3.0)
          {$\mathrm{C_3H_6}\!\mid\!\mathrm{C_3H_8}$\\[1pt]{\Large\bfseries 99.5 wt\%}};
        \draw[main] (raw) -- (deb);
        \draw[main] (deb) -- (rx);
        \draw[main] (rx) -- (de);
        \draw[main] (de) -- (mc);
        \draw[main] (mc) -- (prod);
        % PSA 枝（上）＋ H2 目標
        \node[sep, minimum width=1.3cm] (psa) at (6.0,5.2) {PSA};
        \draw[br] (de.north) -- node[left, font=\scriptsize]{軽質} (psa.south);
        \node[goalH] (h2) at (8.5,5.2)
          {$\mathrm{H_2}\!\mid\!$軽質\\[1pt]{\Large\bfseries 99.9 mol\%}};
        \draw[br] (psa.east) -- (h2.west);
        % 除去・循環（下、補助）
        \node[sub] (c4) at (2.0,1.35) {$\mathrm{C_4}$ 除去};
        \draw[br, black!55] (deb.south) -- (c4.north);
        \node[sub] (rec) at (8.1,1.35) {$\mathrm{C_3H_8}$ 循環};
        \draw[br, black!55] (mc.south) -- (rec.north);
      \end{tikzpicture}}

      \vspace{0.1em}

      % --- 下：膜有無の比較（レポート Fig7.1 形式のままスライド用に拡大）---
      \centering
      \includegraphics[width=\linewidth]{membrane_vs_dist3_slide}\par
      \vspace{0.08em}
      \raggedright
      {\scriptsize 膜無しは塔単独で $200$ 段・塔高 $156$ m を要する。}
    \end{column}
  \end{columns}

\end{frame}
